\begin{table}[H]
\centering
\caption{Migration as a Mechanism: Agricultural Intensity and the Second Great Migration}
\label{tab:migration}
\begin{threeparttable}
\begin{tabular}{lcc}
\toprule
 & (1) & (2) \\
 & Migrated 1940--50 & Occ.\ Gain 1930--50 \\
 & (All Farm Workers) & (Black Farm Workers) \\
\midrule
Farm Share $\times$ Black & 0.0188$$ & \\
 & (0.0143) & \\
Farm Share & 0.0331$^{***}$ & \\
 & (0.0115) & \\
Black & 0.0007$$ & \\
 & (0.0091) & \\
Farm Share $\times$ Migrant & & 2.827$^{***}$ \\
 & & (0.477) \\
Farm Share & & -3.381$^{***}$ \\
 & & (0.303) \\
Migrant & & 4.102$^{***}$ \\
 & & (0.337) \\
\midrule
State FE & & Yes \\
Observations & 2,241,483 & 388,508 \\
Clusters & 1,086 & 982 \\
\bottomrule
\end{tabular}
\begin{tablenotes}[flushleft]
\small
\item \textit{Notes:} Standard errors clustered at the county level in parentheses. * $p<0.10$, ** $p<0.05$, *** $p<0.01$. Column (1): dependent variable is an indicator for interstate migration between 1940 and 1950; sample includes all farm workers (Black and white). Column (2): dependent variable is the 20-year change in occupational score (1930 to 1950); sample restricted to Black farm workers only. ``Migrant'' indicates interstate move between 1940 and 1950.
\end{tablenotes}
\end{threeparttable}
\end{table}

